\documentclass[11pt]{amsart}
\usepackage{amssymb}
\usepackage{mathtools}
\usepackage[expansion=false]{microtype}
\usepackage[letterpaper,margin=1in]{geometry}
\usepackage{enumitem}
\usepackage[hidelinks]{hyperref}
\usepackage{url}
\urlstyle{tt}
\usepackage[capitalise,noabbrev]{cleveref}

\setlist{itemsep=1pt,topsep=3pt}
\renewcommand{\labelitemi}{$\bullet$}
% Take the section sign from the Computer Modern symbol font.
\DeclareTextSymbolDefault{\textsection}{OMS}

\newtheorem{theorem}{Theorem}[section]
\newtheorem{lemma}[theorem]{Lemma}
\newtheorem{corollary}[theorem]{Corollary}
\theoremstyle{definition}
\newtheorem{definition}[theorem]{Definition}
\theoremstyle{remark}
\newtheorem{remark}[theorem]{Remark}

\newcommand{\U}{\mathcal U}
\newcommand{\V}{\mathcal V}
\newcommand{\W}{\mathcal W}
\newcommand{\Uz}{\mathcal U_0}
\newcommand{\B}{\mathbf 2}
\newcommand{\one}{\mathbf 1}
\newcommand{\R}{\mathbb R}
\newcommand{\Om}{\Omega}
\newcommand{\crisp}{\mathrel{::}}
\newcommand{\cu}{\varepsilon}
\newcommand{\fl}[1]{{#1}^{\flat}}
\newcommand{\trunc}[1]{\lVert #1\rVert}
\newcommand{\EM}{\mathsf{EM}}
\newcommand{\DNE}{\mathsf{DNE}}
\newcommand{\CEM}{\mathsf{CEM}}
\newcommand{\PE}{\mathsf{PE}}
\newcommand{\const}{\mathsf{const}}
\newcommand{\CrispRep}{\mathsf{CrispRep}}
\newcommand{\Conn}{\mathsf{Conn}}
\newcommand{\Pone}{\mathsf P_1}
\newcommand{\Ptwo}{\mathsf P_2}
\newcommand{\ap}{\mathrel{\#}}
\newcommand{\apP}[1]{\mathrel{\#_{#1}}}
\newcommand{\ag}[1]{\texttt{#1}}

\begin{document}

\title[Tight apartness and excluded middle in cohesive type theory]{Tight
apartness and excluded middle\\ in cohesive type theory}

\author{Claude Opus 5.5}
\author{GPT-6 Astra}
\thanks{Both authors are AI models and contributed equally to the
mathematics. The solution idea, to use the construction of de~Jong and
Swan, came from GPT-6 Astra; Claude Opus 5.5 derived the results from
it, formalized them, and wrote this paper from a draft by GPT-6 Astra.
We thank Colin Zwanziger, who prompted the work and maintains the
repository \url{https://github.com/zwanzigerc/cohesion}, for
posing the questions and coordinating our work.}

\date{26 October 2026}

\begin{abstract}
Two questions on the nLab ask, assuming crisp excluded middle and
Shulman's Axiom~C1 of cohesive type theory, whether every tight apartness
relation on a discrete type is crisp, and whether the relation that a
tight apartness relation on a type $A$ induces on $\flat A$ is decidable.
The answer to both questions is no in real cohesion and yes in classical
trivial cohesion. Under the stated assumptions, each principle is
equivalent to excluded middle in the cohesive mode. Thus, granting
Shulman's intended models of real cohesion, the stated axioms neither
prove nor refute either principle.
\end{abstract}

\maketitle

\section{Introduction}

Shulman's spatial and real-cohesive type theory~\cite{Shulman2018}, which
internalizes Lawvere's axiomatic cohesion~\cite{Lawvere2007}, extends
homotopy type theory by \emph{crisp} variables, which range over a
discrete base, and a modality $\flat$ whose coalgebras are the
\emph{discrete} types. His Axioms C0, C1 and C2 ask that discreteness be
detected by a family of generators, and Axiom R$\flat$ that the real line
suffice. Since topological models refute excluded middle, Shulman assumes
it only for crisp propositions~\cite[\S2.1]{Shulman2018}. An
\emph{apartness relation} is a proposition-valued relation $x\ap y$ that
is irreflexive, symmetric and cotransitive, that is, $x\ap y$ implies
$(x\ap z)\vee(y\ap z)$; it is \emph{tight} if $\neg(x\ap y)$ implies
$x=y$~\cite{TroelstravanDalen1988}. Classically every tight apartness
relation is $\neq$, but constructively a type may carry several.

The nLab page \emph{open problems in homotopy type
theory}~\cite{nLabOpenProblems} listed two questions about these notions,
posed by the nLab contributor Anonymouse. On 24 October 2026 they read as
follows.
\begin{enumerate}[label=(\arabic*)]
\item Suppose a discrete cohesive type $A$ comes with a tight apartness
relation. Then, assuming crisp excluded middle and axiom C1, is the tight
apartness relation on $A$ a crisp relation?
\item Suppose a cohesive type $A$ comes with a tight apartness relation.
Then, assuming crisp excluded middle and axiom C1, is the tight apartness
relation on $\flat A$ in the cohesive mode decidable?
\end{enumerate}
We take Shulman's formulation of Axiom~C1~\cite[\S8.2]{Shulman2018} as
authoritative; \Cref{rem:nlab} discusses the formulation on the nLab.

For a crisp apartness relation both questions are trivial
(\Cref{rem:reading}), so we read them as principles $\Pone$ and $\Ptwo$
about crisp carriers and arbitrary tight apartness relations
(\Cref{def:principles}). We show that $\Ptwo$ is equivalent to excluded
middle $\EM$, and so is $\Pone$ under crisp excluded middle $\CEM$
(\Cref{thm:em}). Under Axiom C1, $\CEM$ and propositional
extensionality, both are equivalent to every generator being
$0$-connected (\Cref{thm:generators}). A type with two distinct points
and only constant maps to the Booleans refutes $\Ptwo$, and $\Pone$ under
$\CEM$ (\Cref{thm:refutation}); in particular Axiom C2 and real cohesion
do (\Cref{cor:C2}). Since the degenerate interpretation with trivial
cohesion validates $\EM$, the axioms $\CEM$ and C1 do not refute the
principles, and, granting Shulman's intended topos models of real
cohesion, they do not prove them (\Cref{rem:independence}). In this
sense both questions have a negative answer.

Everything rests on a construction from work of de~Jong and
Swan~\cite{deJong2025,TypeTopology}. For a proposition $P$, the relation
$x\apP{P}y\coloneqq P\wedge(x\neq y)$ is an apartness relation on any type
with decidable equality, tight when $\neg\neg P$, and on the Booleans
$0\apP{P}1$ holds exactly when $P$ does. A principle that makes tight
apartness relations on the Booleans crisp, or decidable at crisp points,
therefore transfers crisp information to an arbitrary $\neg\neg$-dense
proposition, which gives double negation elimination. De~Jong, following
a generalization of an idea of Swan, used this relation to show that a
type with decidable equality and two distinct points has at most one
tight apartness relation only if double negation elimination
holds~\cite{TypeTopology} (see \Cref{rem:uniqueness}); that work does
not concern cohesion. Applying
the construction to the present questions was the idea of GPT-6 Astra.

The lemmas, theorems and corollaries of Sections~\ref{sec:prelim}
to~\ref{sec:refutations}, apart from the clauses about Axiom R$\flat$ and
the unit generator, are formalized in Agda, on top of the TypeTopology
library~\cite{TypeTopology}, in the repository~\cite{Formalization},
whose \href{https://github.com/zwanzigerc/cohesion/blob/171714333d2c480a273e3b37293a4a7fc2655526/apartness/README.md\#results-and-agda-names}{apartness README} lists the correspondence.

\section{Preliminaries and the two principles}\label{sec:prelim}

We work in Shulman's spatial type theory~\cite[\S\S2--6]{Shulman2018},
with function extensionality, assumed crisply, and propositional
truncations $\trunc{-}$; univalence is not needed. Shulman also assumes
propositional resizing and propositional extensionality
throughout~\cite[\S1]{Shulman2018}; we avoid the former and state each
use of the latter. Crisp variables are written $u\crisp A$, and closed
terms are crisp. For a crisp type $A$, the type $\flat A$ has elements
$\fl u$ for crisp $u\crisp A$, and \emph{$\flat$-induction} says that to
inhabit a family $C$ over $\flat A$, which need not be crisp, at
$M:\flat A$, it suffices to inhabit $C(\fl u)$ for a crisp variable
$u\crisp A$~\cite[\S4.2]{Shulman2018}. The counit $\cu_A:\flat A\to A$
sends $\fl u$ to $u$. A crisp type $A$ is \emph{discrete} if $\cu_A$ is
an equivalence~\cite[Definition~6.12]{Shulman2018}, and \emph{crisply}
discrete if this has a crisp proof. The Booleans $\B=\one+\one$ are
crisply discrete~\cite[Theorems~6.19 and~6.21]{Shulman2018}, a fact that
rests on crisp case analysis~\cite[Theorem~5.4]{Shulman2018}.

All universes are crisp, and $\Uz$ is the smallest. For a universe
$\V$, write $\Om_\V$ for the type of propositions in $\V$, and $\EM_\V$
and $\DNE_\V$ for excluded middle and double negation elimination for
them, which are equivalent. $\CEM_\V$ says
that every crisp proposition $P\crisp\Om_\V$ is decidable, which is
Shulman's Axiom LEM~\cite[\S2.1]{Shulman2018}, and $\PE_\V$ is
propositional extensionality. A statement that quantifies over a crisp
variable, such as $\CEM_\V$, becomes a type by quantifying over an
element of $\flat\Om_\V$, or of $\flat\U$, and using
$\flat$-induction~\cite[Remarks~6.28 and~8.1]{Shulman2018}: for instance,
$\CEM_\V$ becomes $\prod_{B:\flat\Om_\V}(\cu(B)+\neg\cu(B))$. We use this
form whenever such a statement is a hypothesis or is negated. A type has
\emph{decidable equality} if $(x=y)\vee(x\neq y)$ for all $x$ and~$y$.

\subsection*{Axioms of cohesion}
A type $A$ is \emph{$B$-null} if $\const:A\to(B\to A)$ is an
equivalence. Axiom C0~\cite[\S8.1]{Shulman2018} asks for a crisp family
$R\crisp J\to\mathsf{Type}$ such that a crisp type is discrete if and
only if it is $R_j$-null for all $j:J$. Axiom
C1~\cite[\S8.2]{Shulman2018} adds points $p:\prod_{j:J}R_j$, and Axiom
C2~\cite[\S9]{Shulman2018} adds a $j_0:J$ such that $R_{j_0}$ is a set
with elements $r_0\neq r_1$. Axiom R$\flat$~\cite[\S8]{Shulman2018} is
Axiom C0 for the single generator $\R$, the Dedekind reals, and it implies
C2, since $0\neq1$ in $\R$~\cite[\S9]{Shulman2018}. Our arguments never
use that $J$, $R$ or $p$ is crisp. We therefore call any family
$R:J\to\W$ with points $p_j:R_j$ that satisfies the condition of C0 a
family of \emph{pointed generators}, and we take that condition to apply
to crisp types in all universes.

\begin{lemma}[{cf.~\cite[Lemma~8.7]{Shulman2018}}]\label{lem:null}
Let $b:B$. A type $A$ is $B$-null if and only if $f(x)=f(b)$ for all
$f:B\to A$ and $x:B$.
\end{lemma}

\begin{proof}
Evaluation at $b$ is a retraction of $\const$, and by function
extensionality the condition says exactly that it is also a section.
\end{proof}

\begin{remark}[The formulation on the nLab]\label{rem:nlab}
The page~\cite{nLabPunctual} to which ``axiom C1'' in the questions
links uses one pointed type $I$ and two modes: a crisp type $T$ is
discrete if and only if every function $I\to\mathrm{Disc}(T)$ is
constant. Read literally, this applies discreteness, a property of
cohesive-mode types, to the crisp-mode type $T$, while $\mathrm{Disc}(T)$
is always discrete; the page also writes the canonical map as $A\to\flat A$
instead of $\flat A\to A$. The evident meaning is Shulman's Axiom C1
with one generator, which is how we read the questions; ``in the
cohesive mode'' then adds nothing.
\end{remark}

\begin{lemma}[{de~Jong and Escard\'o~\cite{TypeTopology}}]\label{lem:classical}
Under $\DNE_\V$, every $\V$-valued tight apartness relation on a type
satisfies $x\ap y\leftrightarrow x\neq y$.
\end{lemma}

\begin{proof}
By irreflexivity, and since tightness turns $x\neq y$ into
$\neg\neg(x\ap y)$.
\end{proof}

\begin{lemma}[{de~Jong and Swan~\cite{deJong2025,TypeTopology}}]\label{lem:dejong}
Let $X$ have decidable equality and let $P$ be a proposition with
$\neg\neg P$. Then $x\apP{P}y\coloneqq P\times(x\neq y)$ is a tight
apartness relation on $X$, and on $\B$ we have
$0\apP{P}1\leftrightarrow P$.
\end{lemma}

\begin{proof}
For cotransitivity, given $p:P$, $x\neq y$ and $z$, decide whether
$x=z$: if so, $y\apP{P}z$, and otherwise $x\apP{P}z$. For tightness,
assume $\neg(x\apP{P}y)$; then $x\neq y$ would give
$\neg P$, so $\neg\neg(x=y)$, and hence $x=y$ by decidable equality. The
other properties are immediate.
\end{proof}

\begin{definition}\label{def:principles}
Let $X\crisp\U$ be crisp. A relation $R:X\to X\to\V$, not necessarily
crisp, is \emph{crisply representable} if
\[
 \CrispRep(R)\coloneqq
 \Bigl\lVert \sum_{r:\flat(X\to X\to\Om_{\U\sqcup\V})}\
 \prod_{x,y:X}\bigl(\cu(r)(x,y)\leftrightarrow R(x,y)\bigr)\Bigr\rVert .
\]
For crisp universes $\U$ and $\V$, the principles are:
\begin{description}[font=\normalfont,leftmargin=1.5em,labelindent=0pt]
\item[$\Pone(\U,\V)$] for every crisp $X\crisp\U$ with a crisp proof
that $X$ is discrete, every $\V$-valued tight apartness relation on $X$
is crisply representable;
\item[$\Ptwo(\U,\V)$] for every crisp $X\crisp\U$, every $\V$-valued
tight apartness relation $\ap$ on $X$, and all $u,v:\flat X$, the
proposition $\cu(u)\ap\cu(v)$ is decidable.
\end{description}
The apartness relation is an ordinary variable, not assumed crisp.
\end{definition}

\begin{remark}[On the reading]\label{rem:reading}
(a) For a crisp apartness relation, with a crisp proof that its values are
propositions, both questions are trivial: the relation represents itself,
after lifting its values to $\U\sqcup\V$; and for $u,v:\flat X$ we may
assume by $\flat$-induction that $u=\fl s$ and $v=\fl t$ with $s,t\crisp
X$, so that $\cu(u)\ap\cu(v)$ is the crisp proposition $s\ap t$, which
$\CEM$ decides.
(b) Under propositional extensionality, which Shulman assumes,
$\CrispRep(R)$ says that $R$, with values lifted to $\U\sqcup\V$, lies
merely in the image of the counit
$\flat(X\to X\to\Om_{\U\sqcup\V})\to(X\to X\to\Om_{\U\sqcup\V})$, that
is, that $R$ comes from a crisp relation. Pointwise equivalence
expresses this without propositional extensionality.
(c) Under Axiom C1, a crisp proof of discreteness is no
restriction~\cite[Lemma~8.10]{Shulman2018}.
(d) We read the relation on $\flat A$ as the restriction along $\cu_A$.
It is again a tight apartness relation when $\cu_A$ is injective, as for
crisp sets under Axiom C1~\cite[Theorem~8.21]{Shulman2018}.
(e) Excluded middle implies both principles for all universes; we prove
the converses for carriers in $\Uz$, which contains~$\B$.
\end{remark}

\section{Both principles are excluded middle}\label{sec:em}

\begin{theorem}\label{thm:em}
Let $\V$ be a crisp universe.
\begin{enumerate}[label=(\roman*)]
\item $\EM_\V$ implies $\Ptwo(\U,\V)$ for every $\U$, and $\Ptwo(\Uz,\V)$
implies $\EM_\V$.
\item $\EM_\V$ implies $\Pone(\U,\V)$ for every $\U$. If $\CEM_\V$ holds
and $\B$ is crisply discrete, then $\Pone(\Uz,\V)$ implies $\EM_\V$.
\end{enumerate}
Neither part uses propositional extensionality or an axiom of cohesion.
\end{theorem}

\begin{proof}
Assume $\EM_\V$. It decides the values of every apartness relation, which
gives $\Ptwo(\U,\V)$. For $\Pone(\U,\V)$, let $\ap$ be a $\V$-valued tight
apartness relation on a crisp $X\crisp\U$. By \Cref{lem:classical},
$x\ap y\leftrightarrow x\neq y$. The relation
$(x,y)\mapsto\bigl((x=y)\to\mathbf 0_\V\bigr)$, together with the proof
that its values are propositions, involves only the crisp variable $X$
and the crisp function extensionality, so it gives an element $r$ of
$\flat(X\to X\to\Om_{\U\sqcup\V})$ with $\cu(r)(x,y)\leftrightarrow
x\neq y\leftrightarrow x\ap y$.

For the converses we prove $\DNE_\V$. Let $P:\Om_\V$ with $\neg\neg P$. By
\Cref{lem:dejong}, $\apP{P}$ is a $\V$-valued tight apartness relation on
the crisp type $\B:\Uz$. Under $\Ptwo(\Uz,\V)$, the proposition
$\cu(\fl 0)\apP{P}\cu(\fl 1)$, which is $0\apP{P}1$, is decidable; a
positive answer gives $P$, and a negative one gives $\neg P$, which is
impossible. Under $\Pone(\Uz,\V)$, with $\CEM_\V$ and a crisp proof that
$\B$ is discrete, we get $\CrispRep(\apP{P})$. As the goal $P$ is a
proposition, we may assume $r:\flat(\B\to\B\to\Om_\V)$ with equivalences
$\cu(r)(x,y)\leftrightarrow x\apP{P}y$, which need not be crisp. By
$\flat$-induction, $r=\fl s$ with $s$ crisp, so $s(0,1)$ is a crisp
proposition, which $\CEM_\V$ decides. If it holds, then $0\apP{P}1$, and
so $P$; otherwise $\neg(0\apP{P}1)$, which gives $\neg P$, which is
impossible.
\end{proof}

\begin{corollary}\label{cor:em}
Under $\CEM_\V$, the statements $\Pone(\Uz,\V)$, $\Ptwo(\Uz,\V)$ and
$\EM_\V$ are equivalent.
\end{corollary}

\begin{remark}[Relation to de~Jong's theorem]\label{rem:uniqueness}
Under $\CEM_\V$, a crisply representable tight apartness relation on $\B$
is pointwise equivalent to $\neq$: at $(0,1)$ and $(1,0)$, $\CEM_\V$
decides the crisp representative, and the negative case would give $0=1$
by tightness. So, with propositional extensionality, $\Pone(\Uz,\V)$
makes $\neq$ the only $\V$-valued tight apartness relation on $\B$, and
for $\V=\Uz$ de~Jong's uniqueness theorem~\cite{TypeTopology} also gives
$\DNE_{\Uz}$. What is new is this passage from crispness to uniqueness;
the proof of \Cref{thm:em} avoids propositional extensionality.
\end{remark}

\section{Axiom C1 and the generators}\label{sec:generators}

In this section $R:J\to\W$ is a family of pointed generators and $\B$ is
crisply discrete, as under Axiom C1. The generators are
\emph{$0$-connected} if $\Conn(R,p)\coloneqq\prod_{j:J}\prod_{x:R_j}
\trunc{x=p_j}$; this is a condition on identity types, not on the shape
modality.

\begin{lemma}\label{lem:omega}
If the counit $\flat\Om_\V\to\Om_\V$ has a section, for instance if
$\Om_\V$ is discrete, then $\CEM_\V$ implies $\EM_\V$.
\end{lemma}

\begin{proof}
For $P:\Om_\V$, the section gives $Q:\flat\Om_\V$ with $\cu(Q)=P$. By
$\flat$-induction $Q=\fl S$ with $S\crisp\Om_\V$, so $\CEM_\V$ decides
$S$, which is $\cu(Q)$, and hence~$P$.
\end{proof}

\begin{theorem}\label{thm:generators}
\leavevmode
\begin{enumerate}[label=(\roman*)]
\item $\EM_\W$ implies $\Conn(R,p)$.
\item Under $\PE_\W$ and $\CEM_\W$, $\Conn(R,p)$ implies $\EM_\W$.
\end{enumerate}
Hence, under $\PE_\W$ and $\CEM_\W$, the statements $\Pone(\Uz,\W)$,
$\Ptwo(\Uz,\W)$, $\EM_\W$ and $\Conn(R,p)$ are equivalent.
\end{theorem}

\begin{proof}
(i) Since $\B$ is discrete, it is $R_j$-null, so every $f:R_j\to\B$
satisfies $f(x)=f(p_j)$ by \Cref{lem:null}. Apply this to the map
$\chi:R_j\to\B$ that $\EM_\W$ provides, with $\chi(y)=1$ if
$\trunc{y=p_j}$ and $\chi(y)=0$ otherwise: as $\chi(p_j)=1$, we get
$\trunc{x=p_j}$ for every $x:R_j$.
(ii) By $\PE_\W$, the type $\Om_\W$ is a set. So for $f:R_j\to\Om_\W$
and $x:R_j$ the goal $f(x)=f(p_j)$ is a proposition, and we may obtain it
from $\trunc{x=p_j}$. By \Cref{lem:null}, $\Om_\W$ is $R_j$-null for
every $j$, hence discrete by the condition of C0, applied to the crisp
type $\Om_\W$ in the next universe; no resizing is involved.
\Cref{lem:omega} gives $\EM_\W$. The last statement follows with
\Cref{cor:em}.
\end{proof}

\begin{corollary}\label{cor:single}
Let $(I,p)$ be the only generator, with $I:\W$ a set, and assume
$\CEM_\W$. Then $\Pone(\Uz,\W)$, $\Ptwo(\Uz,\W)$, $\EM_\W$ and
$\prod_{x:I}(x=p)$ are equivalent, without propositional extensionality.
In particular, if the only generator is the unit type, $\CEM_\W$ implies
$\EM_\W$.
\end{corollary}

\begin{proof}
By (i), $\EM_\W$ gives $\trunc{x=p}$, which is $x=p$ as $I$ is a set.
Conversely, if $x=p$ for all $x$, then $f(x)=f(p)$ for every
$f:I\to\Om_\W$ without $\Om_\W$ being a set, and we conclude as in
(ii).
\end{proof}

\section{Refutations and independence}\label{sec:refutations}

\begin{theorem}\label{thm:refutation}
Let $I:\W$ have elements $q_0\neq q_1$ such that $f(x)=f(y)$ for every
$f:I\to\B$ and $x,y:I$. Then $\neg\EM_\W$ and $\neg\Ptwo(\Uz,\W)$, and,
if $\CEM_\W$ holds and $\B$ is crisply discrete, also
$\neg\Pone(\Uz,\W)$. This needs no axiom of cohesion, no assumption that
$I$ is a set, and no propositional extensionality.
\end{theorem}

\begin{proof}
Under $\EM_\W$, let $\chi:I\to\B$ be $1$ exactly where $\trunc{x=q_0}$
holds. Then $\chi(q_1)=\chi(q_0)=1$, so $\trunc{q_1=q_0}$, contradicting
$q_0\neq q_1$. Now apply \Cref{thm:em}.
\end{proof}

\begin{corollary}\label{cor:C2}
If some pointed generator $R_{j_0}:\W$ has elements $r_0\neq r_1$, as in
Axiom C2, and $\B$ is crisply discrete, then $\neg\EM_\W$ and
$\neg\Ptwo(\Uz,\W)$, and under $\CEM_\W$ also $\neg\Pone(\Uz,\W)$. In
particular, Axiom R$\flat$ refutes $\Ptwo(\Uz,\W)$, and together with
Axiom LEM also $\Pone(\Uz,\W)$, for every universe $\W$ containing~$\R$.
\end{corollary}

\begin{proof}
Every $f:R_{j_0}\to\B$ satisfies $f(x)=f(p_{j_0})=f(y)$, as in the proof
of \Cref{thm:generators}, so \Cref{thm:refutation} applies. Axiom
R$\flat$ implies C2 with the generator $\R$~\cite[\S9]{Shulman2018}.
\end{proof}

\begin{remark}[An explicit family]\label{rem:family}
Excluded middle is known to fail in real
cohesion~\cite[\S8.3]{Shulman2018}. Assume Axiom R$\flat$, and for $x:\R$
let $Q(x)$ be $(x>0)\vee\neg(x>0)$, a proposition in $\Uz$ for cuts
valued in $\Om_{\Uz}$. As $\neg\neg Q(x)$, \Cref{lem:dejong} gives a family
$x\mapsto{\apP{Q(x)}}$ of $\Uz$-valued tight apartness relations on $\B$.
If all of them were decidable at crisp points, or crisply representable
under Axiom LEM, the proof of \Cref{thm:em} would give $Q(x)$ for all
$x:\R$, making $x>0$ a detachable subset of $\R$. By
\cite[Theorem~11.1]{Shulman2018} it would then be empty or all of $\R$,
which is absurd. So $\Pone(\Uz,\Uz)$ and $\Ptwo(\Uz,\Uz)$ fail too, as
witnessed by apartness relations depending on a real parameter.
\end{remark}

\begin{remark}[Independence]\label{rem:independence}
In the interpretation where $\flat$ is the identity and every type is
discrete, which spatial type theory allows~\cite[\S7]{Shulman2018},
Axiom C1 holds with the unit type as generator, and over a classical
metatheory $\EM$ holds, so both principles do. Hence $\CEM$ and C1 do
not refute them. Conversely, Shulman explains that Axiom LEM holds in a
local topos over a base satisfying excluded middle, that Axiom C2
expresses punctual local connectedness and contractible codiscreteness,
and that Axiom R$\flat$ describes his motivating model, sheaves on the
Cartesian spaces $\R^n$~\cite[\S\S7 and~10]{Shulman2018}. In such models
both principles fail by \Cref{cor:C2}, so $\CEM$ and C1 do not prove
them. Shulman calls this semantics
conjectural~\cite[Remark~7.5]{Shulman2018}; we do not formalize these
remarks.
\end{remark}

\begin{thebibliography}{9}
\setlength{\itemsep}{0pt}

\bibitem{deJong2025}
T.~de~Jong.
\newblock A little thread on constructive mathematics and (tight)
  apartness relations in particular.
\newblock Mathstodon, February 2025.
\newblock \url{https://mathstodon.xyz/@de_Jong_Tom/113930642479158387}.

\bibitem{Formalization}
\emph{Tight apartness relations in cohesive type theory}.
\newblock Agda formalization of the results of this paper, 2026.
\newblock \url{https://github.com/zwanzigerc/cohesion}.

\bibitem{Lawvere2007}
F.~W. Lawvere.
\newblock Axiomatic cohesion.
\newblock \emph{Theory and Applications of Categories}, 19(3):41--49,
  2007.

\bibitem{nLabOpenProblems}
nLab authors.
\newblock Open problems in homotopy type theory.
\newblock
  \url{https://ncatlab.org/nlab/show/open+problems+in+homotopy+type+theory}.
\newblock Consulted 24 and 26 October 2026.

\bibitem{nLabPunctual}
nLab authors.
\newblock Axiom of punctual cohesion.
\newblock
  \url{https://ncatlab.org/nlab/show/axiom+of+punctual+cohesion}.
\newblock Revision of 14 May 2025.

\bibitem{Shulman2018}
M.~Shulman.
\newblock Brouwer's fixed-point theorem in real-cohesive homotopy type
  theory.
\newblock \emph{Mathematical Structures in Computer Science},
  28(6):856--941, 2018.
\newblock Numbering as in
  \href{https://arxiv.org/abs/1509.07584v3}{arXiv:1509.07584v3}.

\bibitem{TroelstravanDalen1988}
A.~S. Troelstra and D.~van Dalen.
\newblock \emph{Constructivism in Mathematics: An Introduction}, volumes~I
  and~II.
\newblock North-Holland, 1988.

\bibitem{TypeTopology}
M.~H. Escard\'o and contributors.
\newblock TypeTopology, \url{https://github.com/martinescardo/TypeTopology},
  commit \ag{52ed2eb}.
\newblock Module \ag{Apartness.Properties}, by T.~de~Jong and
  M.~H. Escard\'o.

\end{thebibliography}

\end{document}
